\section{Materials and methods}

\subsection{Data curation}

Each of the 609 dossiers originates from a MIBiG accession and its associated primary reference(s).
Structured fields (BGC identity and boundaries, organism and strain, genes, metabolites, experiments,
measurements, claims, relationships, and source citations) were extracted from the primary literature
and the MIBiG record using Claude (Anthropic) large language models
\todo{Methods -- state the exact Claude model version(s) and extraction date range once confirmed; left as a placeholder rather than guessed}.
Every one of the 609 dossiers was then \textbf{fully manually reviewed} by the curator against the
cited literature before being included in this snapshot; this is a completeness property of the
current release, not a partial-sample audit.

Extraction outputs were passed through a normalisation layer (the \texttt{TERM} vocabulary and
\texttt{NORMALIZATION\_QUEUE} tables) that reconciles free-text labels for experimental conditions,
interventions, assays and taxonomy against a controlled vocabulary where possible. Values that could
not yet be confidently normalised (composite labels, type mismatches, or terms awaiting review) are
retained rather than forced into a canonical bucket, and are tracked explicitly: 882 \texttt{TERM}
rows remained in the normalisation queue at the time of this snapshot (Results). Similarly, where
independent sources disagreed, or a boundary or causal claim was ambiguous, this was recorded as a
\texttt{CONFLICT} row with an explicit issue type and status rather than resolved by editorial
fiat -- 995 such records exist in the current snapshot, of which 220 remain explicitly \emph{open}.

\subsection{Database construction}

All code is released under \texttt{scripts/} in the project repository and is orchestrated by a
single entry point, \texttt{run\_all.py}, so the entire resource is deterministically rebuildable
from the source workbook:

\begin{enumerate}
  \item \textbf{\texttt{parse\_xlsx.py}} reads the 23 data tables and 5 documentation sheets from the
    source workbook and writes one clean CSV per table.
  \item \textbf{\texttt{build\_database.py}} loads the parsed tables into a documented SQLite
    database (\texttt{data/db/fungal\_bgc\_atlas.db}, schema in \texttt{db/schema.sql}), with every
    table's row count asserted against the source workbook's own summary counts as a build-time
    regression check.
  \item \textbf{\texttt{build\_graph.py}} constructs two graphs: a literal knowledge graph of the
    curated subject--predicate--object \texttt{RELATIONSHIP} triples, and an entity-linkage graph
    joining BGC, organism, gene, metabolite and source nodes on \texttt{dossier\_id}, exported as
    GraphML, JSON and CSV edge/node lists.
  \item \textbf{\texttt{compute\_stats.py}} computes every descriptive statistic used in this
    manuscript and in the dashboard directly from the parsed tables, and emits both CSV and
    \LaTeX{}-ready table snippets (\texttt{tables/}).
  \item \textbf{\texttt{make\_figures.py}} renders dated, publication-quality figure snapshots
    (300\,dpi PNG + vector SVG) to \texttt{figures/<date>/}, with an accompanying manifest.
  \item \textbf{\texttt{build\_dashboard.py}} renders the interactive dashboard
    (\texttt{docs/index.html}) described in Results.
\end{enumerate}

The relational schema comprises 23 tables sharing \texttt{dossier\_id} as a common join key (Table~1
of the online Data Dictionary panel; full schema in \texttt{db/schema.sql}), spanning identity and
administration (\texttt{BGC}, \texttt{ORGANISM}, \texttt{STRAIN}), independent categorisation
(\texttt{BGC\_FACET}, 27 axes), molecular content (\texttt{GENE}, \texttt{METABOLITE}), experimental
context (\texttt{EXPERIMENT}, \texttt{TERM}, \texttt{EXPERIMENT\_LINK}), results
(\texttt{MEASUREMENT}, \texttt{OBSERVATION}), interpretive records (\texttt{CLAIM},
\texttt{RELATIONSHIP}, \texttt{CONFLICT}), provenance (\texttt{SOURCE}, \texttt{EVIDENCE\_LINK},
\texttt{ASSERTION}), and curation-audit tables (\texttt{COMPACT\_FACT}, \texttt{COVERAGE},
\texttt{DOSSIER\_SCHEMA}, \texttt{NORMALIZATION\_QUEUE}, \texttt{PROVENANCE\_EXCLUDED},
\texttt{RAW\_ROW}).

The interactive dashboard is a single self-contained HTML page using Plotly for charting and a
small vanilla-JavaScript component for the searchable dossier browser, requiring no server or
build step to view. A GitHub Actions workflow rebuilds the database, graph, statistics, figures and
dashboard on every push to the main branch and deploys the result to GitHub Pages, so the published
resource is always reproducible from the checked-in source data and code.
